% GENERATED FILE. Edit canonical lessons, then run npm run generate:books.
% canonical-source-hash: fnv1a64:c940906e39a9af52
% canonical-lessons: JA-C469-gakki, JA-C469-suzu, JA-C469-tataku, JA-C469-kanaderu, JA-C469-neiro

\chapter{Instruments}
\label{ch:ja-instruments}
\hlchaptermodality{469}

\begin{chapteropening}
\textbf{By the end of this chapter:} \emph{I can talk about musical instruments.}

Complete the last lesson of chapter 469: I can talk about musical instruments.
\end{chapteropening}

\section[\texorpdfstring{\ja{がっき} (gakki)}{がっき (gakki)}]{\ja{がっき} (gakki) — a musical instrument}

\label{lesson:JA-C469-gakki}



\begin{quote}
Before the new one: say the Japanese for a masterpiece, then the Japanese for to assemble.
\end{quote}

\subsection*{You'll want to know: \ja{がっき}}
\textbf{\ja{がっき}} — \emph{gakki} — \textquotedblleft{}a musical instrument\textquotedblright{}.

\begin{quote}
\textbf{\ja{あねは} \ja{がっきが} \ja{じょうずです。}} — \emph{ane wa gakki ga jōzu desu} — \textquotedblleft{}My elder sister is good at instruments.\textquotedblright{}
\end{quote}

\begin{cousinweb}
\textbf{\ja{がっき}} (\emph{gakki}) is Sino-Japanese (kango): two Chinese characters read with their on'yomi, readings borrowed from Chinese. They are \emph{ga}, \textquotedblleft{}music\textquotedblright{} (Middle Chinese \emph{ngaewk}) and \emph{ki}, \textquotedblleft{}instrument\textquotedblright{} (\emph{khijH}) (KANJIDIC2; Middle Chinese in Baxter's transcription).
\end{cousinweb}

\begin{grammarlens}[title={What you've built}]
One more word for describing what you lived and felt.
\end{grammarlens}

\subsection*{Your turn}
\begin{itemize}
\raggedright
  \item \emph{Say it:} \emph{gakki}
  \item \emph{Say it:} \emph{gakki}, once more
  \item \emph{Say it:} \emph{ane wa gakki ga jōzu desu} — My elder sister is good at instruments.
  \item \emph{Recall:} say the Japanese for a masterpiece, then the Japanese for to assemble, then say \emph{gakki} again
\end{itemize}

\begin{tcolorbox}[breakable,colback=teal!4,colframe=teal!35!black,title={Before you move on}]
What does \ja{がっき} mean? (\textquotedblleft{}A musical instrument\textquotedblright{}.) Say it once more.
\end{tcolorbox}
\section[\texorpdfstring{\ja{すず} (suzu)}{すず (suzu)}]{\ja{すず} (suzu) — a small bell}

\label{lesson:JA-C469-suzu}



\begin{quote}
Before the new one: say the Japanese for to assemble, then the Japanese for a musical instrument.
\end{quote}

\subsection*{You'll want to know: \ja{すず}}
\textbf{\ja{すず}} — \emph{suzu} — \textquotedblleft{}a small bell\textquotedblright{}.

\begin{quote}
\textbf{\ja{ねこの} \ja{くびに} \ja{すずが} \ja{あります。}} — \emph{neko no kubi ni suzu ga arimasu} — \textquotedblleft{}The cat has a little bell on its neck.\textquotedblright{}
\end{quote}

\begin{cousinweb}
\textbf{\ja{すず}} (\emph{suzu}) is a native Japanese word (wago), inherited rather than borrowed from Chinese; its kanji spelling was attached to it later (JMdict).
\end{cousinweb}

\begin{grammarlens}[title={What you've built}]
One more word for describing what you lived and felt.
\end{grammarlens}

\subsection*{Your turn}
\begin{itemize}
\raggedright
  \item \emph{Say it:} \emph{suzu}
  \item \emph{Say it:} \emph{suzu}, once more
  \item \emph{Say it:} \emph{neko no kubi ni suzu ga arimasu} — The cat has a little bell on its neck.
  \item \emph{Recall:} say the Japanese for to assemble, then the Japanese for a musical instrument, then say \emph{suzu} again
\end{itemize}

\begin{tcolorbox}[breakable,colback=teal!4,colframe=teal!35!black,title={Before you move on}]
What does \ja{すず} mean? (\textquotedblleft{}A small bell\textquotedblright{}.) Say it once more.
\end{tcolorbox}
\section[\texorpdfstring{\ja{たたく} (tataku)}{たたく (tataku)}]{\ja{たたく} (tataku) — to beat; to tap}

\label{lesson:JA-C469-tataku}



\begin{quote}
Before the new one: say the Japanese for a musical instrument, then the Japanese for a small bell.
\end{quote}

\subsection*{You'll want to know: \ja{たたく}}
\textbf{\ja{たたく}} — \emph{tataku} — \textquotedblleft{}to beat; to tap\textquotedblright{}.

\begin{quote}
\textbf{\ja{こどもが} \ja{たいこを} \ja{たたいています。}} — \emph{kodomo ga taiko o tataite imasu} — \textquotedblleft{}The children are beating a drum.\textquotedblright{}
\end{quote}

Polite \textbf{\ja{たたきます}}, \emph{tatakimasu}; a godan verb. Past (polite): \textbf{\ja{おとうとが} \ja{たいこを} \ja{たたきました。}}, \emph{otōto ga taiko o tatakimashita}.

\begin{cousinweb}
\textbf{\ja{たたく}} (\emph{tataku}) is a native Japanese word (wago), inherited rather than borrowed from Chinese; its kanji spelling was attached to it later (JMdict).
\end{cousinweb}

\begin{grammarlens}[title={What you've built}]
One more word for describing what you lived and felt.
\end{grammarlens}

\subsection*{Your turn}
\begin{itemize}
\raggedright
  \item \emph{Say it:} \emph{tataku}
  \item \emph{Say it:} \emph{tataku}, once more
  \item \emph{Say it:} \emph{kodomo ga taiko o tataite imasu} — The children are beating a drum.
  \item \emph{Recall:} say the Japanese for a musical instrument, then the Japanese for a small bell, then say \emph{tataku} again
\end{itemize}

\begin{tcolorbox}[breakable,colback=teal!4,colframe=teal!35!black,title={Before you move on}]
What does \ja{たたく} mean? (\textquotedblleft{}To beat; to tap\textquotedblright{}.) Say it once more.
\end{tcolorbox}
\section[\texorpdfstring{\ja{かなでる} (kanaderu)}{かなでる (kanaderu)}]{\ja{かなでる} (kanaderu) — to play music}

\label{lesson:JA-C469-kanaderu}



\begin{quote}
Before the new one: say the Japanese for a small bell, then the Japanese for to beat a drum.
\end{quote}

\subsection*{You'll want to know: \ja{かなでる}}
\textbf{\ja{かなでる}} — \emph{kanaderu} — \textquotedblleft{}to play music\textquotedblright{}.

\begin{quote}
\textbf{\ja{かのじょは} \ja{ふえで} \ja{うたを} \ja{かなでました。}} — \emph{kanojo wa fue de uta o kanademashita} — \textquotedblleft{}She played a song on the flute.\textquotedblright{}
\end{quote}

Polite \textbf{\ja{かなでます}}, \emph{kanademasu}; an ichidan verb. Past (polite): \textbf{\ja{せんせいが} \ja{ふえを} \ja{かなでました。}}, \emph{sensei ga fue o kanademashita}.

\begin{cousinweb}
\textbf{\ja{かなでる}} (\emph{kanaderu}) is a native Japanese word (wago), inherited rather than borrowed from Chinese; its kanji spelling was attached to it later (JMdict).
\end{cousinweb}

\begin{grammarlens}[title={What you've built}]
One more word for describing what you lived and felt.
\end{grammarlens}

\subsection*{Your turn}
\begin{itemize}
\raggedright
  \item \emph{Say it:} \emph{kanaderu}
  \item \emph{Say it:} \emph{kanaderu}, once more
  \item \emph{Say it:} \emph{kanojo wa fue de uta o kanademashita} — She played a song on the flute.
  \item \emph{Recall:} say the Japanese for a small bell, then the Japanese for to beat a drum, then say \emph{kanaderu} again
\end{itemize}

\begin{tcolorbox}[breakable,colback=teal!4,colframe=teal!35!black,title={Before you move on}]
What does \ja{かなでる} mean? (\textquotedblleft{}To play music\textquotedblright{}.) Say it once more.
\end{tcolorbox}
\section[\texorpdfstring{\ja{ねいろ} (neiro)}{ねいろ (neiro)}]{\ja{ねいろ} (neiro) — the tone of an instrument}

\label{lesson:JA-C469-neiro}



\begin{quote}
Before the new one: say the Japanese for to beat a drum, then the Japanese for to play music.
\end{quote}

\subsection*{You'll want to know: \ja{ねいろ}}
\textbf{\ja{ねいろ}} — \emph{neiro} — \textquotedblleft{}the tone of an instrument\textquotedblright{}.

\begin{quote}
\textbf{\ja{この} \ja{ふえは} \ja{ねいろが} \ja{きれいです。}} — \emph{kono fue wa neiro ga kirei desu} — \textquotedblleft{}This flute has a beautiful tone.\textquotedblright{}
\end{quote}

\begin{cousinweb}
\textbf{\ja{ねいろ}} (\emph{neiro}) is a native Japanese compound (wago), inherited rather than borrowed from Chinese: \emph{ne}, \textquotedblleft{}sound\textquotedblright{} and \emph{iro}, \textquotedblleft{}colour\textquotedblright{} (JMdict, KANJIDIC2).
\end{cousinweb}

\begin{grammarlens}[title={What you've built}]
One more word for describing what you lived and felt.
\end{grammarlens}

\subsection*{Your turn}
\begin{itemize}
\raggedright
  \item \emph{Say it:} \emph{neiro}
  \item \emph{Say it:} \emph{neiro}, once more
  \item \emph{Say it:} \emph{kono fue wa neiro ga kirei desu} — This flute has a beautiful tone.
  \item \emph{Recall:} say the Japanese for to beat a drum, then the Japanese for to play music, then say \emph{neiro} again
\end{itemize}

\begin{tcolorbox}[breakable,colback=teal!4,colframe=teal!35!black,title={Before you move on}]
What does \ja{ねいろ} mean? (\textquotedblleft{}The tone of an instrument\textquotedblright{}.) Say it once more.
\end{tcolorbox}
